\scanpage{13}
- 194 - \hfill n\textsuperscript{o} 391

5. \uline{Systèmes de type} $A_\ell$, $(\ell\geq1)$.

(I) Considérons dans l'espace $V=\mathrm R^{\ell+1}$ l'intersection du réseau $L_{\unsure{o|0}}$, (n\textsuperscript{o} 4), et de l'hyperplan $V'$ d'équation $\sum_{i=1}^{\ell+1}\xi_i=0$ ; c'est un réseau $L'_{\unsure{o|0}}$ dans $V$ ; la valeur minima de $(u|u)$ dans $L'_{\unsure{o|0}}$ est $2$ et elle est atteinte pour les vecteurs de la forme $\varepsilon_i-\varepsilon_j$, $(i\neq j,\ 1\leq i\leq\ell+1,\ 1\leq j\leq\ell+1)$. Montrons que l'ensemble $R$ de ces vecteurs est un \uline{système de racines}, (évidemment réduit), dans $V'$. En effet, il est clair que pour $u$, $v$ dans $R$, $2(u|v)/(u|u)$ est entier puisque $(u|u)=2$ et la conclusion résulte du n\textsuperscript{o} 4. Le nombre des racines $n$ est $\ell(\ell+1)$.

(II) Posons :
\[
\alpha_1=\varepsilon_1-\varepsilon_2\ ,\quad
\alpha_2=\varepsilon_2-\varepsilon_3\ ,\ldots,\quad
\alpha_\ell=\varepsilon_\ell-\varepsilon_{\ell+1}\ .
\]
Il est immédiat que tout élément de $R$ est combinaison linéaire des $\alpha_i$, $(1\leq i\leq\ell)$, à coefficients tous $\geq0$ ou tous $\leq0$ ; donc, (§ 6, n\textsuperscript{o} 7, cor. 3 de la prop. 14), les $\alpha_i$ forment un système de \uline{racines simples} $S$ de $R$. En outre, on a $\|\alpha_i\|^2=2$ pour tout $i$, $(\alpha_i|\alpha_{i+1})=-1$ pour $1\leq i\leq\ell-1$ et $(\alpha_i|\alpha_j)=0$ pour $j>i+1$ ; le schéma de Dynkin de $R$ est donc du type $A_\ell$, ce qui montre que $R$ est irréductible.

(III) En vertu du § 8, n\textsuperscript{o} 3, formule (4), on a $h=n/\ell=\ell+1$.

(IV) L'ordre de la plus grande racine est $h-1=\ell$, (§ 8, n\textsuperscript{o} 5, prop. 6). Or, $\varepsilon_1-\varepsilon_{\ell+1}=\alpha_1+\alpha_2+..+\alpha_\ell$ est une racine positive d'ordre $\ell$ ; donc c'est la plus grande racine $\widetilde\alpha$, (§ 6, n\textsuperscript{o} 6, prop. 13). Elle est orthogonale à toutes les racines simples sauf à $\alpha_1$ et $\alpha_\ell$, et on a $(\widetilde\alpha|\alpha_1)=(\widetilde\alpha|\alpha_\ell)=1$.
\mspage{12}
